\subsection{等价函数类}

设 \(\mu\) 是局部紧空间 \(X\) 上的一个测度。给定一个集合 \(F\) ，两个映射 \(f, g\) （从 \(X\) 到 \(F\) 中）称为关于 \(\mu\) 等价的（或 \(\mu\) -等价的，或在不致混淆时简称等价的），如果 \(f(x) = g(x)\) 在 \(X\) 中几乎处处成立。由于两个可忽略集之并是可忽略的，这样确实在集 \(F^X\) —— 由 \(X\) 到 \(F\) 中的全体映射组成 —— 中定义了一个等价关系；当我们说到这样一个函数 \(f\) 的等价类（不加其他说明）时，指的将是与 \(f\) 几乎处处相等的全体函数所成的类；在本章及以后各章中，我们将用记号 \(\widetilde{f}\) 表示这个类。

命题 8. --- 设 \((\mathrm{F}_{n})\) 是集合的一个可数族（有限或无穷）。对每个指标 n，设 \(f_{n}, g_{n}\) 是从 X 到 \(F_{n}\) 中的两个等价映射；则存在一个可忽略集 H ，使得对每个 \(x \notin H\) ，对一切 n 有 \(f_{n}(x) = g_{n}(x)\) 。

事实上，集 \(H_{n}\) —— 由诸 \(x \in X\) （满足 \(f_{n}(x) \neq g_{n}(x)\) ）组成 —— 是可忽略的，因此它们的并 H 也可忽略（No.~2, 命题 4），而这个集满足要求。

推论. --- 若 \(\varphi\) 是从 \(\prod_{n} F_n\) 到集合 \(G\) 中的一个映射，则映射 \(\varphi((f_n))\) 与 \(\varphi((g_n))\) （从 \(X\) 到 \(G\) 中）是等价的。

我们用 \(\varphi((\widetilde{f}_n))\) 表示每个函数 \(\varphi((f_n))\) 的等价类，其中 \(f_n\) 是类 \(\widetilde{f}_n\) 中的任意函数。

特别地，若 F 是 R 上的向量空间，则定义 \(\widetilde{f} + \widetilde{g}\) 与 \(\alpha\widetilde{f}\) 分别为 \(f + g\) 与 \(\alpha f\) 的等价类（f 与 g 是从 X 到 F 中的映射，而 \(\alpha\) 为标量）；这样，我们在从 X 到 F 中的映射的等价类所成之集上得到一个向量空间结构：而且，这就是 \(F^{X}\) 的结构关于由满足 \(\widetilde{f} = \widetilde{0}\) 的映射 f 组成的线性子空间（即几乎处处为零的函数，我们也称之为（取值于 F 的）可忽略函数）的商空间结构。类似地定义乘积 \(\widetilde{g}\widetilde{f}\) ，其中 \(\widetilde{f}\) 是从 X 到 F 中的映射的一个等价类，而 \(\widetilde{g}\) 是定义在 X 上的（有限）数值函数的一个等价类：这样，从 X 到 F 中的映射的等价类所成之集就具有一个在定义于 X 上的有限数值函数的等价类所成之集（它自身具有一个环结构）上的模结构。若 F 是 R 上的代数，则类似地在从 X 到 F 中的映射的等价类所成之集上定义一个代数结构。

设 F 是可度量的拓扑空间，并考虑 F 上一个与其拓扑相容且由伪距离的可数族 \(\rho_{n}\) 定义的一致结构（GT, IX, §§1 与 2）；要使从 X 到 F 中的两个映射 f, g 等价，必须且只须数值函数 \(\rho_{n}(f,g)\) 都可忽略；事实上，这等价于说 X 中存在一个可忽略集 H ，使得对每个 \(x\notin H\) ，对一切 n 有 \(\rho_{n}(f(x),g(x))=0\) ，即 \(f(x)=g(x)\) 。特别地，若 F 是可度量的局部凸空间，且 \((q_{n})\) 是定义 F 的拓扑的可数半范数族（TVS, II, §4, No.~1），则要使从 X 到 F 中的两个映射 f, g 等价，必须且只须数值函数 \(q_{n}(\mathbf{f}(x)-\mathbf{g}(x))\) 都可忽略。

命题 9. --- 设 f 与 g 是从 X 到 Hausdorff 拓扑空间 F 中的两个连续映射；要使 f 与 g 等价，必须且只须 \(f(x) = g(x)\) 在 \(\mu\) 的支集的每点处成立。

事实上，由诸 \(x \in X\) （满足 \(f(x) \neq g(x)\) ）组成的集是一个开集（GT, I, §8, No.~1）；要使它可忽略，必须且只须它与 \(\mu\) 的支集不相交（No.~2, 命题 5）。

命题 10. --- 设 F 是 R 上的 Hausdorff 局部凸空间，使得在 F 的对偶 F' 中存在一个序列 ( \(a_{n}^{\prime}\) ) ，它在弱拓扑 \(\sigma(F^{\prime},F)\) 下稠密（TVS, II, §6, No.~2）。要使从 X 到 F 中的两个映射 f, g 等价，必须且只须对每个 n，数值函数 \(\langle\mathbf{f}(x),\mathbf{a}_{n}^{\prime}\rangle\) 与 \(\langle\mathbf{g}(x),\mathbf{a}_{n}^{\prime}\rangle\) 等价。

该条件显然是必要的。反之，若它被满足，则存在一个可忽略集 H ，使得对每个 \(x \notin H\) ，对一切 n 有 \(\langle \mathbf{f}(x), \mathbf{a}_{n}^{\prime} \rangle = \langle \mathbf{g}(x), \mathbf{a}_{n}^{\prime} \rangle\) ；这表明弱连续线性形式 \(\mathbf{z}^{\prime} \mapsto \langle \mathbf{f}(x), \mathbf{z}^{\prime} \rangle\) 与 \(\mathbf{z}^{\prime} \mapsto \langle \mathbf{g}(x), \mathbf{z}^{\prime} \rangle\) （在 \(F^{\prime}\) 上）在各点 \(a_{n}^{\prime}\) 处相等，从而由假设它们恒等，这就证明了 \(\mathbf{f}(x) = \mathbf{g}(x)\) 对所有 \(x \notin H\) 成立。

注意，命题 10 的假设特别适用于下述情形：F 是可度量且可分的 \(^{1}\) 局部凸空间（TVS, III, §3, No.~4, 命题 6 的推论 2）。

